\documentclass[a4paper, 12pt]{article}

\usepackage[utf8]{inputenc}
\usepackage[T1]{fontenc}
\usepackage{textcomp}
\usepackage{amssymb}
\usepackage{newtxtext} \usepackage{newtxmath}
\usepackage{amsmath, amssymb}
\newtheorem{problem}{Problem}
\newtheorem{example}{Example}
\newtheorem{lemma}{Lemma}
\newtheorem{theorem}{Theorem}
\newtheorem{problem}{Problem}
\newtheorem{example}{Example} \newtheorem{definition}{Definition}
\newtheorem{lemma}{Lemma}
\newtheorem{theorem}{Theorem}
\DeclareMathAlphabet{\mathcal}{OMS}{cmsy}{m}{n}

\begin{document}


Given an EEG signal, we define SWA activity (following Guo et al.) as the
amplitude of the combined signal resulting from the IMFs with frequencies in the
range of 0.2-4 Hz. 

Guo et al. model transition from SWS to not-SWS as a stochastic process of
forming a synchronized neural cluster of size $S(t)$, with a trivial model: 

\begin{equation}
    S(t+1) = S(t) + \Delta S
\end{equation}

where $\Delta S$ is a random variable. They then make the model more complex,
observing that the random variable $\Delta S$ should be proportional to the
cluster size, i.e. $\Delta S \prop S(t)$, producing 

\begin{equation}
    S(t+1) = S(t) + \delta S(t) = (1+\alpha) S(t)
\end{equation}

Taking the logarithm, 

\begin{equation}
    \log S(t+1) = \log S(t) + \log(1+\delta) 
\end{equation}

Here, $\log(1 + \delta)$ is a new random variable. This is a random walk model. 
This implies that the distribution of $S(t)$ is log-normal, and that the
distribution of $\log S(t)$ is normal.

SWA is taken as a proxy measure of the cluster size $S(t)$. The authors then
show that the distribution 




\end{document}



